\documentclass[10pt,a4paper]{article}

\usepackage[margin=0.65in]{geometry}
\usepackage[T1]{fontenc}
\usepackage{lmodern}
\usepackage{enumitem}
\usepackage[hidelinks]{hyperref}
\usepackage{titlesec}
\usepackage{xcolor}
\usepackage{qrcode}

%====================
% Color Theme
%====================
\definecolor{PrimaryColor}{HTML}{1F4E79}
\definecolor{SecondaryColor}{HTML}{4B5563}

\setlength{\parindent}{0pt}
\setlist[itemize]{leftmargin=*, noitemsep, topsep=2pt}

\titleformat{\section}
  {\large\bfseries\color{PrimaryColor}}
  {}
  {0em}
  {}
  [\color{SecondaryColor}\titlerule]

\hypersetup{
    colorlinks=true,
    urlcolor=PrimaryColor,
    linkcolor=PrimaryColor
}

\urlstyle{same}

\begin{document}

%====================
% Header with QR
%====================
\begin{minipage}[t]{0.78\textwidth}
    {\LARGE \textbf{\color{PrimaryColor}Abdelrahman Elsaid Elsawy}}\\
    {\color{SecondaryColor}Robotik ve Araç Simülasyonu Mühendisi}\\
    ROS/ROS2 \textbar{} BeamNG.tech \textbar{} NVIDIA Isaac Sim \textbar{} RL Eğitimi \textbar{} ONNX \textbar{} Otomatik QA\\
    E-posta: \href{mailto:abdelrahman.elsaid.elsawy@gmail.com}{abdelrahman.elsaid.elsawy@gmail.com}\\
    LinkedIn: \href{https://www.linkedin.com/in/abdelrahman-elsaid-elsawy}{linkedin.com/in/abdelrahman-elsaid-elsawy}\\
    Telefon/WhatsApp: \href{https://wa.me/905365231899}{+90 536 523 1899}\\
    Konum: İstanbul, Türkiye
\end{minipage}
\hfill
\begin{minipage}[t]{0.18\textwidth}
    \raggedleft
    \qrcode[height=1.7cm]{https://www.linkedin.com/in/abdelrahman-elsaid-elsawy}\\
    {\scriptsize LinkedIn QR}
\end{minipage}

%====================
\section*{Özet}

ROS/ROS2, sistem entegrasyonu, otonom sistemler ve yüksek sadakatli araç simülasyonu alanlarında 6 yılı aşkın deneyime sahip robotik ve yazılım mühendisi. BeamNG.tech, ROS/ROS2, MATLAB ve Simulink dahil simülasyon platformlarını robotik ve mühendislik araç zincirlerine bağlama konusunda uzmanlaşmıştır. API köprüsü geliştirme, gerçek zamanlı simülasyon kontrolü, veri toplama, otomatik QA testi, regresyon analizi, sorun araştırma, teknik destek ve B2B simülasyon iş akışları için müşteriye dönük mühendislik konularında deneyimlidir. Ek odak alanları arasında NVIDIA Isaac Sim, pekiştirmeli öğrenme (RL) eğitim iş akışları, ONNX model dağıtımı ve MATLAB kullanılarak Yapay Sinir Ağları (YSA) tabanlı araştırma analizi yer almaktadır.

%====================
\section*{Temel Yetkinlikler}

\textbf{\color{PrimaryColor}Robotik:} ROS, ROS2, SLAM, MoveIt, Gazebo, RViz, otonom navigasyon, sensör füzyonu\\
\textbf{\color{PrimaryColor}Simülasyon:} BeamNG.tech, NVIDIA Isaac Sim, araç simülasyonu, yüksek sadakatli simülasyon, simülasyon iş akışları\\
\textbf{\color{PrimaryColor}Yapay Zeka ve Model Dağıtımı:} Pekiştirmeli öğrenme, RL eğitim iş akışları, ONNX, yapay sinir ağları, YSA tabanlı analiz\\
\textbf{\color{PrimaryColor}Programlama:} C, C++, Python, Lua, SQL, Bash, XML, JSON\\
\textbf{\color{PrimaryColor}Entegrasyon:} API köprüleri, sistem entegrasyonu, gerçek zamanlı simülasyon kontrolü, veri toplama, MATLAB, Simulink\\
\textbf{\color{PrimaryColor}Test ve QA:} otomatik QA testi, regresyon testi, hata yeniden üretimi, kök neden analizi, sorun araştırma\\
\textbf{\color{PrimaryColor}Donanım:} Arduino, ESP32, Raspberry Pi, LiDAR, IMU, ultrasonik sensörler, SPI/I2C/UART, PCB tasarımı, havya lehimleme, prototipleme, Catia\\
\textbf{\color{PrimaryColor}Araçlar:} Git, Docker, Ubuntu, Zendesk, API dokümantasyonu

%====================
\section*{Profesyonel Deneyim}

\textbf{\color{PrimaryColor}Research Software Engineer} \hfill Kas 2021 -- Günümüz\\
\textbf{BeamNG GmbH} \hfill Uzaktan

\begin{itemize}
    \item BeamNG.tech platformunu ROS/ROS2, MATLAB ve Simulink ile yüksek sadakatli araç simülasyonu, robotik iş akışları ve otonom sistem geliştirme için entegre etti.
    \item Gerçek zamanlı simülasyon kontrolü, araç etkileşimi, veri toplama ve harici araç entegrasyonu için API köprüleri geliştirdi ve bakımını üstlendi.
    \item Motor kararlılığı, simülasyon güvenilirliği ve entegrasyon kalitesini iyileştirmek için C++, Python ve Lua kullanarak otomatik QA testleri, regresyon analizi ve sorun araştırması gerçekleştirdi.
    \item Müşteri kaynaklı ve dahili olarak tespit edilen teknik sorunları; hataları yeniden üreterek, API davranışlarını inceleyerek, senaryo yapılarını analiz ederek, kök nedenleri belirleyerek araştırdı; düzeltmeler, geçici çözümler ve uygulama rehberliği sağladı.
    \item Zendesk üzerinden simülasyon, API, ROS/ROS2, MATLAB/Simulink ve entegrasyon sorunlarını çözerek küresel B2B müşterilere destek verdi.
    \item BeamNG.tech API kullanımı, MATLAB/Simulink iş akışları, eğitimler, HIL/DIL materyalleri ve entegrasyon rehberliği için geliştirici dokümantasyonunu iyileştirdi.
    \item Tekrarlayan QA ve müşteri desteği bulgularını, geliştirici deneyimini iyileştirmek için yeniden kullanılabilir test senaryolarına, dokümantasyon iyileştirmelerine ve uygulanabilir ürün geri bildirimlerine dönüştürdü.
    \item BeamNGpy, API dokümantasyonu, simülasyon iş akışları, QA düzeltmeleri, eğitimler ve entegrasyon örnekleri dahil BeamNG.tech mühendislik depolarına katkıda bulundu.
    \item Programatik senaryo keşfi, senaryo yükleme, araç getirme, sensör bağlama, araç bağlantısı, komut dosyasıyla kontrol ve sensör sorgulama dahil senaryo ve görev yükleme iş akışlarını uyguladı ve doğruladı.
\end{itemize}

\textbf{\color{PrimaryColor}Mechatronics Systems Engineer} \hfill Şub 2020 -- Tem 2021\\
\textbf{Ideasparq Robotics Sdn Bhd} \hfill Yerinde

\begin{itemize}
    \item Servis robotları için ROS tabanlı kontrol sistemleri ve otonom navigasyon çözümleri geliştirdi.
    \item Sensör füzyonu algoritmaları uyguladı ve robotik platformları Gazebo, RViz, Docker ve Ubuntu tabanlı robotik iş akışlarıyla test etti.
    \item Yazılım, sensörler, kontrol mantığı, simülasyon ve dağıtım ortamları genelinde robotik sistem entegrasyonu üzerinde çalıştı.
\end{itemize}

\textbf{\color{PrimaryColor}Autonomous Robotics Engineer} \hfill Haz 2019 -- Oca 2020\\
\textbf{ENDSTRUCT} \hfill Yerinde

\begin{itemize}
    \item GOVTECH mobil robotları için otonom navigasyon ve sağlık izleme sistemleri geliştirdi.
    \item Navigasyon, izleme ve sahaya yönelik robotik sistem davranışlarını içeren mobil robotik yazılımları üzerinde çalıştı.
\end{itemize}

\textbf{\color{PrimaryColor}Research Assistant} \hfill Eyl 2019 -- Ara 2019\\
\textbf{Autonomous Vehicle Laboratory}

\begin{itemize}
    \item Gerçek zamanlı nesne tespiti için 3B LiDAR verilerini işledi.
    \item Sensör soyutlaması ve otonom araç araştırma iş akışları için ROS tabanlı ara yazılım geliştirdi.
    \item Mühendislik araştırma analizi ve modele dayalı tahmin iş akışları için yapay sinir ağı kavramlarını uyguladı.
\end{itemize}

\textbf{\color{PrimaryColor}Mekatronik Stajyeri} \hfill Haz 2018 -- Eyl 2018\\
\textbf{University Malaysia Pahang}

\begin{itemize}
    \item TurtleBot3 ve ROS robot simülasyonu ile uygulamalı deneyim kazandı.
\end{itemize}

%====================
\section*{Seçili Teknik Katkılar}

\begin{itemize}
    \item BeamNG.tech'in ROS/ROS2, MATLAB ve Simulink ile entegrasyonu.
    \item Gerçek zamanlı simülasyon kontrolü ve veri erişimi için BeamNGpy ve API köprüsü geliştirme.
    \item Otomatik simülasyon iş akışları için senaryo ve görev yükleme doğrulaması.
    \item QA otomasyonu, regresyon testi, hata yeniden üretimi ve kök neden analizi.
    \item Küresel B2B simülasyon kullanıcıları için müşteriye dönük mühendislik desteği.
    \item API'ler, eğitimler, entegrasyon iş akışları, MATLAB/Simulink ve HIL/DIL materyalleri için dokümantasyon iyileştirmeleri.
    \item ROS, Gazebo, RViz, Docker, Ubuntu ve NVIDIA Isaac Sim kullanarak robotik simülasyon ve entegrasyon.
    \item Robotik ve simülasyon tabanlı otonom sistemler için pekiştirmeli öğrenme eğitim iş akışları bilgisi.
    \item Taşınabilir YZ model değişimi ve çıkarım (inference) iş akışları için ONNX model dağıtımı bilgisi.
    \item MATLAB odaklı iş akışlarını kullanarak YSA tabanlı mühendislik araştırma analizi.
\end{itemize}

%====================
\section*{Eğitim}

\textbf{\color{PrimaryColor}Mekatronik, Robotik ve Otomasyon Mühendisliği Lisansı} \hfill 2016 -- 2020\\
UMPSA Universiti Malaysia Pahang Al-Sultan Abdullah

%====================
\section*{Yayın}

\textbf{\color{PrimaryColor}The Flexural Strength Prediction of Porous Cu-Sn-Ti Composites via Artificial Neural Networks}\\
The 6th International Conference on Robotics Intelligence and Applications, 2018\\
DOI: \href{https://link.springer.com/chapter/10.1007/978-981-13-8323-6_34}{10.1007/978-981-13-8323-6\_34}

%====================
\section*{Kurslar \& Sertifikalar}

\begin{itemize}[noitemsep, topsep=5pt]
    \item C1 İngilizce Seviyesi -- Busuu (Mayıs 2025)
    \item A1 Almanca Seviyesi -- Busuu (Kas 2024)
    \item SQL ile İlişkisel Veritabanları -- Manara (Kas 2024)
    \item Google Sheets ile Temel Veri Analizi -- Coursera (Eki 2024)
    \item Veritabanı Temelleri -- Manara (Eki 2024)
    \item Lua Scripting: Sıfırdan Tam Lua Programlama -- Udemy (Kas 2023)
    \item Temel Veri Koruma Eğitimi (AB) -- Mitarbeiterschule GmbH (Eyl 2023)
    \item ROS1 Geliştiricisi Olarak ROS2 Öğrenme ve ROS Projelerini Taşıma -- Udemy (Kas 2022)
    \item Seviye 2: SLAM için Simülasyon ve Donanım Uygulaması -- Udemy (Kas 2022)
    \item Azure DevOps Boards'a Giriş -- Coursera (Eki 2022)
    \item Linux'ta Komut Satırı -- Coursera (Eyl 2022)
    \item Arduino Atölyesi -- Core Electronics (Eyl 2022)
    \item Açık Kaynak Yazılım Geliştirme Yöntemleri -- The Linux Foundation (Eki 2021)
    \item DevOps için Python Scripting Uzmanlığı -- LearnQuest (Eyl 2021)
    \item Python Scripting: Dosyalar, Kalıtım ve Veritabanları -- LearnQuest (Eyl 2021)
    \item Web Uygulaması Teknolojileri ve Django -- University of Michigan (Eyl 2021)
    \item AWS Cloud Technical Essentials -- AWS (Ağu 2021)
    \item Python ile DevOps ve Derleme Otomasyonu -- LearnQuest (Ağu 2021)
    \item Python Scripting: Tarihler, Sınıflar ve Koleksiyonlar -- LearnQuest (Ağu 2021)
    \item Yapılandırma Yönetimi ve Bulut -- Google (Ağu 2021)
    \item Git ve GitHub'a Giriş -- Google (Tem 2021)
    \item Teknik Destek Temelleri -- Google (Tem 2021)
    \item Sorun Giderme ve Hata Ayıklama Teknikleri -- Google (Tem 2021)
    \item İşletim Sistemiyle Etkileşim İçin Python Kullanımı -- Google (Haz 2021)
    \item Python Hızlı Kurs -- Google (May 2021)
    \item Malezya Otonom Araç Yarışması (MAVeC'19) -- MARii (Kas 2018)
    \item Robotik, Zeka ve Uygulamalar Uluslararası Konferansı (RiTA) 2018 -- KIAST SG (2018)
    \item Boundless Innovation Week -- HPMM (Ağu 2018)
    \item CAM (Bilgisayar Destekli Üretim) Atölyesi -- Malezya Mühendisler Odası (Kas 2017)
    \item IEM MATLAB Atölyesi -- Malezya Mühendisler Odası (Eki 2017)
    \item Yeni Başlayanlar için Arduino Atölyesi -- SPARKTECH Team at IIUM (Kas 2016)
    \item Çevrimiçi Veritabanı Dersi: SCOPUS -- University Malaysia Pahang (Kas 2016)
\end{itemize}

%====================
\section*{Diller}

Arapça: Anadil veya iki dilli yetkinlik\\
İngilizce: Mesleki çalışma yetkinliği\\
Türkçe: Temel seviye\\
Almanca: Temel seviye\\
Malayca: Sınırlı çalışma yetkinliği

\end{document}